\clearpage
\begin{summarybox}
\textbf{\textcolor{CSummary}{最短记忆：奇函数看“和”，偶函数看“距离”}}

设函数定义在 $\mathbb R$ 或包含 $0$ 的左右对称区间，且在右半边\textbf{严格}单调。$a,b$ 均在定义域内。
\[
\begin{array}{c|c|c}
\text{函数}&\text{右半边严格递增}&\text{右半边严格递减}\\\hline
\text{奇函数}&f(a)+f(b)>0\iff a+b>0&f(a)+f(b)>0\iff a+b<0\\[2pt]
\text{偶函数}&f(a)>f(b)\iff|a|>|b|&f(a)>f(b)\iff|a|<|b|
\end{array}
\]

\textbf{\textcolor{CSummary}{为什么成立？}}
\begin{enumerate}[leftmargin=*,itemsep=0.35em]
  \item \textbf{奇函数：}由 $f(-b)=-f(b)$，把 $f(a)+f(b)>0$ 变成 $f(a)>f(-b)$；奇函数右半边严格递增（或递减）时，整个区间也是同方向严格单调，最后只比 $a$ 与 $-b$。
  \item \textbf{偶函数：}由 $f(x)=f(|x|)$，把 $a,b$ 都搬到 $0$ 右边，再比 $|a|$ 与 $|b|$。
\end{enumerate}

\textbf{\textcolor{CSummary}{两个顺手得到的结论}}
\[
\begin{aligned}
\text{奇函数：}\quad&f(a)+f(b)=0\iff a=-b;\\
\text{偶函数：}\quad&f(a)=f(b)\iff a=b\ \text{或}\ a=-b.
\end{aligned}
\]
奇函数严格递增时 $f(x)>0\iff x>0$；偶函数右半边严格递增时，$f(0)$ 是最小值（严格递减时则是最大值）。

\medskip
\textbf{\textcolor{CSummary}{做题顺序}}

\textbf{1. 看奇偶 $\rightarrow$ 2. 看严格增减 $\rightarrow$ 3. 变成 $a+b$ 或 $|a|,|b|$ $\rightarrow$ 4. 解普通不等式 $\rightarrow$ 5. 检查定义域。}

\textbf{\textcolor{CSummary}{特别提醒：}}只知道“不下降”不能套严格等价；奇函数的两个数求和规则不能随意改成三个数求和规则；偶函数不能直接用 $a>b$ 比较函数值。
\end{summarybox}
